\documentclass[10pt,a4paper]{article}
\usepackage[margin=18mm,top=22mm,bottom=22mm]{geometry}
\usepackage{kotex}
\usepackage{graphicx}
\usepackage{booktabs}
\usepackage{multirow}
\usepackage{array}
\usepackage{amsmath,amssymb}
\usepackage{hyperref}
\hypersetup{hidelinks}
\setlength{\parindent}{10pt}
\setlength{\parskip}{0pt}
\setlength{\columnsep}{9mm}
\begin{document}
\begin{center}
{\LARGE\bfseries 그래프 신경망 기반 다중 환율 예측 : MTGNN 변형과 앙상블 예측 \par}
\vspace{2mm}
{\large\bfseries Graph Neural Network-based Multi-Currency FX Forecasting: MTGNN Variants, Ensemble Prediction \par}
\vspace{3mm}
{\normalsize 김수민 1 ․ 최지원 2 ․ 백승환 3 ․ 한용섭 4 ․ 박대승 5 \par}
{\normalsize Kim Soo-min, Choi Ji-won, Baek Seung-hwan, Han Yong-seop, Park Dae-seung \par}
\vspace{2mm}
{\small 1 수원대학교 컴퓨터학부 SW 학과 \par}
{\small E-mail: [ 본인 이메일 입력 ]@suwon.ac.kr \par}
{\small 2 수원대학교 컴퓨터학부 SW 학과 \par}
{\small E-mail: 23017097 @suwon.ac.kr \par}
{\small 3 수원대학교 컴퓨터학부 SW 학과 \par}
{\small E-mail: sqortmdghks1@suwon.ac.kr \par}
{\small 4 수원대학교 컴퓨터학부 SW 학과 \par}
{\small E-mail: [ 본인 이메일 입력 ]@suwon.ac.kr \par}
{\small 5 수원대학교 컴퓨터학부 SW 학과 \par}
{\small E-mail: [ 본인 이메일 입력 ]@suwon.ac.kr \par}
\end{center}
\vspace{1mm}
\noindent\textbf{요약}\quad 본 연구는 환율 예측의 전통적인 시계열 모델이 간과해 온 국가 간 상호의존성을 포착하는 데 있어 한계를 극복하고 , 환율 변동의 복잡한 메커니즘을 규명하는 데 목적이 있다 . 이를 위해 그래프 신경망 (GNN) 기반의 MTGNN 모델을 채택하여 IMF 주요 3 개국의 월간 환율 데이터와 거시경제 지표를 결합한 시공간적 상관관계를 통합 분석하였다 . 특히 데이터로부터 잠재적 관계를 스스로 학습하는 적응형 그래프 학습 모듈을 통해 예측의 정교함과 확장성을 검증하였다 . 연구 결과 , 제안 모델은 기존 LSTM 및 통계 모델 대비 우수한 예측 정확도를 기록하였다 . 수치 예측의 한계를 넘어 SLM(Small Language Model) 기반의 다중 에이전트 시스템을 결합함으로써 분석의 신뢰성을 강화하였으며 , 전문화된 에이전트 기법으로 순환형 체이닝 구조를 통해 예측값의 시계열 일관성을 검증하였다 . 또한 특정 오차의 원인을 식별하여 모델의 개선 방향을 제시하고 , 장기 예측 결과에 대한 국가별 상호작용 해석과 설명 가능한 (Explainable) 인사이트 도출을 지원한다 . 결론적으로 본 연구는 인공지능 기술을 활용하여 수치 예측과 경제적 맥락 해석이 통합된 차세대 외환 분석 프레임워크를 정립하였다는 데 학술적 의의가 있다 .
\par
\noindent\textbf{키워드 - 환율 예측 , 그래프 신경망 (GNN), 다중 에이전트 시스템 (MAS), 설명 가능한 AI(XAI), 시계열 일관성 검증}
\par
\vspace{1mm}
\noindent\textbf{Abstract}\quad This study aims to overcome the limitations of traditional time-series models in capturing cross-country dependencies and to elucidate the complex mechanisms of exchange rate fluctuations. By adopting the MTGNN (Multivariate Time Series Forecasting with Graph Neural Networks) model, this research conducts an integrated analysis of spatio -temporal correlations by combining monthly exchange rate data from three major IMF countries with macroeconomic indicators. In particular, the precision and scalability of the forecasts are validated through an adaptive graph learning module that autonomously identifies latent relationships within the data. Experimental results demonstrate that the proposed model achieves superior prediction accuracy compared to existing LSTM and statistical models. Beyond numerical forecasting, the reliability of the analysis is further enhanced by integrating a multi-agent system based on a Large Language Model (LLM). This specialized agent technique utilizes a recursive chaining structure to verify the time-series consistency of forecasted values while identifying the causes of specific errors to suggest model improvements. Furthermore, this approach facilitates the interpretation of inter-country interactions and the derivation of explainable insights for long-term forecasts. Consequently, this research holds significant academic value by establishing a next-generation foreign exchange analysis framework that integrates numerical prediction and economic contextual interpretation using artificial intelligence.
\par
\noindent\textbf{Keywords — Exchange Rate Forecasting, Graph Neural Networks (GNN), Multi-Agent System (MAS), Explainable AI (XAI), Time-series Consistency Verification}
\par
\vspace{2mm}
\twocolumn
\section{개요}
최근 글로벌 금융 시장은 국가 간 경제적 밀착도가 심화됨에 따라 환율 변동의 복잡성과 상호의존성이 급격히 증대되고 있다 [1]. 환율은 단순한 개별 국가의 경제 지표를 넘어 , 인접 국가 및 주요 교역국 간의 유기적인 관계 속에서 실시간으로 변화하는 동적 특성을 지닌다 [2]. 그러나 ARIMA 와 같은 전통적인 통계 모델이나 단순 LSTM 계열의 딥러닝 모델들은 이러한 국가 간의 공간적 상관관계를 충분히 반영하지 못하고 , 단일 국가의 과거 패턴 학습에 치중해 왔다는 한계가 있다 [3]. 이는 급변하는 매크로 경제 상황 속에서 국가 간 금융 전이 효과를 간과하게 만들며 , 결과적으로 환율의 장기적 추세를 정확히 포착하고 변동 원인을 해석하는 데 걸림돌이 된다 [4].

이에 본 연구는 환율이 개별적으로 존재하는 것이 아니라 상호 동적인 관계로 연결되어 있다는 점에 주목하여 , 국가 간 잠재적 관계를 스스로 학습하는 그래프 신경망 (Graph Neural Networks, GNN) 기반의 MTGNN 모델을 제안한다 [5]. 특히 명시적인 인접 행렬이 존재하지 않는 외환 시장의 특성을 고려하여 , 적응형 그래프 학습 모듈 (Adaptive Graph Learning Layer) 을 통해 주요 3 개국 ( 미국 , 한국 , 일본 ) 간의 숨겨진 시공간적 정보를 통합 모델링하였다 . 또한 , 금리 및 물가 등 국가별 거시경제 지표를 노드 피처 (Node Feature) 로 결합함으로써 , 모델이 단순 시계열 패턴을 넘어 다차원적인 경제적 맥락을 학습할 수 있도록 확장성을 부여하였다 [6].

본 연구의 핵심적인 차별점은 단순한 수치 예측력 향상에 그치지 않고 , LLM(Large Language Model) 기반의 다중 에이전트 시스템 (Multi-Agent System) 을 결합하여 예측 결과에 대한 고도의 해석력을 확보했다는 점에 있다 [7]. 기존 연구들이 단일 수치 예측에 치중했던 것과 달리 , 본 프레임워크는 전문화된 다중 에이전트 간의 순환형 체이닝 (Chaining) 구조를 도입하여 예측 데이터의 시계열 일관성을 엄밀히 검증하고 심층적인 원인 분석을 수행한다 [8]. 이를 통해 환율 변동의 동인에 대한 설명 가능한 (Explainable) 근거를 구축하며 , 에이전트의 다각적 해석이 담긴 장기 예측 보고서를 제공함으로써 외환과 글로벌 경제에 대한 근본적인 분석을 제안한다 .

단순한 기술적인 정확도 개선을 넘어 , 데이터 기반의 정량적 예측과 경제적 추론 기반의 정성적 분석이 결합된 새로운 외환 분석 패러다임을 제시한다는 점에서 본 연구는 중요한 의의를 갖는다 . 특히 인공지능이 도출한 수치에 전문가적 해석을 덧붙임으로써 , 금융의 실질적인 활용 가능성을 입증하고자 한다 . 이는 변동성이 큰 글로벌 경제 환경에서 보다 선제적이고 입체적인 리스크 대응 체계를 구축하는 데 기여할 것으로 기대된다 .

본 논문의 구성은 다음과 같다 . 2 장에서는 환율 예측 및 거시경제 요인에 관한 선행 연구를 고찰하고 , 3 장에서는 제안하는 MTGNN 및 다중 에이전트 시스템의 상세 구조를 설명한다 . 4 장에서는 실험 과정과 결과를 통해 모델의 유효성을 검증하며 , 마지막으로 결론 및 향후 연구 방향에 대해 논한다 .

\section{관련 연구 HY 중고딕 11pt, 굵게 )}
\subsection{외환 환율 예측 및 딥러닝 모델}
외환 환율 예측 분야는 2000 년대 이전부터 계량경제학적 접근을 통해 활발히 연구되었다 . 초기 연구들은 주로 벡터 자기회귀 (VAR) 나 GARCH 와 같은 통계적 모델을 사용하여 환율의 변동성을 설명하고자 하였다 . [9]. 그러나 이러한 전통적인 방식은 통계적 가정을 필요로 하며 , 금융 시장의 복잡한 비선형 역학을 포착하는 데 한계가 존재한다 [10]. 이러한 한계를 해결하기 위해 최근에는 기계학습 및 딥러닝 기반의 방법론이 제안되었다 .

딥러닝 모델인 LSTM, CNN, 그리고 최근 자연어 처리 분야에서 혁신을 일으킨 Transformer 등이 외환 예측에 적용되어 기존 통계 모델 대비 우수한 성능을 보였다 . 특히 시계열 데이터의 장기 의존성 (Long-term dependency) 과 미시적 변동성을 포착하기 위해 Zhou et al.(2021) 은 시간적 특징을 임베딩 (Temporal Embeddings) 하여 어텐션 메커니즘을 개선한 Informer 모델을 제안하였다 [11]. 이 연구는 단변량 데이터보다 거시경제 지표를 포함한 다변량 (Multivariate) 데이터를 사용할 때 트랜스포머가 LSTM 이나 ARIMA 보다 월등한 예측 성능과 수익률을 달성함을 입증하였다 .

\subsection{거시경제 요인과 기존 방법론의 한계}
선행 연구들은 외환 시장의 예측 정확도를 높이기 위해 다양한 시도를 해왔다 . 금리평가설 (IRP) 이나 구매력평가설과 같은 이론에 기반하여 금리 , 인플레이션율 , 통화 정책과 같은 거시경제 지표를 특징 (feature) 으로 포함하는 연구들이 진행되었다 [12]. Yao et al.(2018) 의 연구 역시 기술적 지표 외에 펀더멘털 ( Fundamental) 데이터를 결합한 Dual-Stage Attention 구조를 통해 예측력을 강화하였다 [13].

그러나 이러한 최신 연구들조차 다음과 같은 구조적 한계를 지닌다 . 첫째 , 다중 통화 간의 구조적 상호 의존성 (Interrelationships) 을 간과하였다 . 대부분의 딥러닝 기반 연구는 입력 데이터를 평면적인 벡터 형태로 처리할 뿐 , 경제학적 이론에 근거한 통화 삼각 관계 (triplet) 나 통화 간의 네트워크 구조를 명시적으로 모델링하지 않는다 . 둘째 , 시장 마찰 요인에 따른 가치 평가의 상대성을 충분히 반영하지 못했다 . 신용도나 은행 규제 등으로 인해 통화의 가치 평가는 상대 통화에 따라 달라질 수 있으나 , 기존 연구는 이를 전체 통화 네트워크 관점에서 포괄적으로 활용하지 못하는 경향이 있다 .

\subsection{그래프 신경망 (GNN) 기반의 다변량 시계열 예측}
앞서 언급한 외환 시장의 구조적 상호 의존성을 포착하기 위해 , 최근 시계열 예측 분야에서는 그래프 신경망 (Graph Neural Networks, GNN) 을 활용한 연구가 주목받고 있다 [6], [14]. 다변량 시계열 데이터에서 각 변수를 그래프의 노드 (Node) 로 간주할 경우 , GNN 은 인접한 노드들의 정보를 집계하여 잠재된 공간적 의존성 (Spatial Dependency) 을 효과적으로 학습할 수 있다는 장점이 있다 [15] ,[16]. 특히 기존의 CNN 이나 RNN 기반 방법론들이 변수 쌍 (Pair-wise) 간의 관계를 명시적으로 모델링하지 못해 모델의 해석력이 떨어진다는 한계를 GNN 을 통해 극복할 수 있다 [17].

그러나 외환 시장과 같은 일반적인 다변량 시계열 데이터는 변수 간의 연결 구조 (Graph Structure) 가 사전에 정의되어 있지 않다는 문제점이 존재한다 . 이를 해결하기 위해 Wu et al.(2020) 은 사전에 정의된 그래프 구조 없이도 데이터로부터 변수 간의 관계를 자동으로 추출할 수 있는 그래프 학습 모듈 (Graph Learning Module) 을 제안하였다 [5]. 이 연구에서 제안한 MTGNN(Multivariate Time Series Grpah Neural Network) 프레임워크는 그래프 학습 레이어를 통해 변수 간의 단방향 (Uni-directed) 관계를 학습하여 인접 행렬 (Adjacency Matrix) 을 생성하고 , 이를 바탕으로 그래프 컨볼루션 (Graph Convolution) 을 수행하여 공간적 관계를 포착한다 .

또한 , 시계열 데이터의 시간적 패턴을 학습하기 위해 팽창 컨볼루션 (Dilated Convolution) 과 인셉션 레이어 (Inception Layer) 를 결합한 시간 컨볼루션 모듈을 함께 사용하여 , 시공간적 (Spatial-Temporal) 의존성을 동시에 모델링하는 종단간 (End-to-End) 학습 구조를 확립하였다 [5]. 이러한 접근 방식은 환율과 같이 변수 간의 구조적 관계가 뚜렷하지 않은 데이터에서도 잠재된 상호 연관성을 발견하고 예측 성능을 향상시킬 수 있음을 시사한다 .

\section{Methodology}
\subsection{문제 정의 및 개요}
본 연구는 다수 국가 ( 또는 통화 ) 로 구성된 글로벌 외환 (FX) 시스템에서 , 국가 간 상호의존성 ( 그래프 구조 ) 과 시계열 동학 ( 시간 의존성 ) 을 함께 고려하여 다중 스텝 (Multi-horizon) 환율 경로를 확률적으로 예측하는 문제를 다룬다 .

국가 ( 또는 통화 ) 집합을 라 하고 , 각 국가는 FX 네트워크의 노드로 표현한다 . 관측 시점 동안 국가별 환율 및 거시 · 금융 특징이 주어졌을 때 , 목표는 다음의 조건부 예측분포를 추정하는 것이다 .

여기서 는 예측 구간 (horizon) 이며 , 는 미래 스텝의 환율 ( 또는 환율 상태 / 수익률 ) 벡터 시퀀스이다 . 이 정의는 한 점 예측이 아니라 , 미래 경로 전체에 대한 분포를 추정한다 .

예를 들어 학습 구간이 2011-01 부터 T=2025-10 까지라면 , T+1 은 2025-11 이고 , H=36 이면 2025-11 부터 36 개월 구간을 예측한다 . 이때 모델은 와 그래프 에 조건부로 의 확률분포를 추정한다 .

\subsection{입력 변수 및 특징 구성}
환율은 단일 요인으로 움직이지 않으며 , ( i ) 환율 자체의 동학 , (ii) 변동성 · 유동성 ( 시장 활동 ), (iii) 거시 · 정책 환경 ( 공통 컨텍스트 ) 의 결합에 의해 비대칭적으로 공진화 (co-evolve) 한다 . 이에 따라 본 연구는 각 국가 i 에 대해 다음의 구성요소를 결합한다 .

\subsubsection{환율 동학}
국가 i 의 환율 관측을 라 할 때 , 로그수익률 ( 또는 정규화된 환율 상태 ) 을 로 정의한다 .

실험 / 데이터 특성에 따라 를 정규화된 레벨로 두고 변환을 생략할 수도 있으나 , 본문에서는 를 환율 상태 ( 수익률 포함 ) 로 정의한다 .

\subsubsection{변동성 · 시장 활동}
환율 리스크를 반영하기 위해 롤링 변동성 와 , 유동성 / 거래강도를 대변하는 거래량 ( 또는 대체 지표 ) 를 포함한다 .

\subsubsection{거시 · 정책 컨텍스트}
여러 국가에 동시에 영향을 주는 거시 · 금융 공통요인을 로 둔다 . ( 예 : 글로벌 위험선호 , 금리 , 원자재 등 )

\subsubsection{최종 특징 벡터}
최종적으로 국가 i 의 시점 t 입력 특징을 다음과 같이 정의한다 .

시장 휴일 / 비동기 캘린더 / 거시지표 결측 등으로 발생하는 결측은 전진 채움 (forward filling) 및 마스킹으로 처리한다 .

\subsection{국가 상호작용 그래프 구성}
국가 간 구조적 의존성을 명시적으로 반영하기 위해 가중 그래프를 정의한다 .

여기서 는 국가 ( 통화 ) 노드 집합 , 는 국가 간 관계 , 는 가중 인접행렬이다 . 엣지 는 무역노출 , 금융통합 , 통화정책 스필오버 , 지역 인접성 , 환율 상관 기반 동조화 등으로 인해 i 의 변화가 j 에 영향을 줄 수 있음을 의미한다 .

초기 인접행렬 는 상관 , 경제 / 지역 근접성 , 잠재 의존 성분을 혼합해 구성할 수 있으며 ,

처럼 가중 결합 형태로 표현된다 . 학습 과정에서 그래프 구조는 메시지 패싱 등을 통해 동적으로 보정 (refinement) 될 수 있다 .

\subsection{베이지안 MTGNN 예측}
\subsubsection{모델 입력 / 출력 및 예측 목표}
시점 t 에서 국가별 특징 행렬을

로 두고 , 각 행이 에 대응한다 . 관측 및 그래프 가 주어졌을 때 , 목적은 미래 환율 상태의 예측분포를 추정하는 것이다 .

\subsubsection{베이지안 예측 ( 불확실성 포함 ) 과 구간 ( 밴드 ) 해석}
추론 시 드롭아웃을 활성화한 채 S 회 확률적 전방추론을 수행한다 .

이로부터 예측 평균과 분산을 다음과 같이 근사한다 .

그러나 MC dropout 은 모델의 인식론적 불확실성 (Epistemic Uncertainty) 만 포착하며 , 데이터의 고유한 노이즈인 우연적 불확실성 (Aleatoric Uncertainty) 을 반영하지 못한다 [18]. 특히 dropout 비율이 낮을 경우 ( ), MC 샘플 간 분산이 극히 작아 예측구간이 지나치게 좁아진다는 한계가 있다 .

이를 해결하기 위해 본 연구는 잔차 기반 예측구간 (Residual-based Prediction Interval) 방법을 제안한다 . 검증 데이터에서 실제 예측 오차의 분산 ( ) 을 추정하고 , 이를 MC dropout 분산과 결합하여 현실적인 예측구간을 산출한다 .

여기서 는 구간 너비 조정 계수로 , 사용자가 시각화 목적에 맞춰 조정할 수 있다 . 본 연구의 실험에서는 를 적용하여 과도한 불확실성 팽창을 억제하였다 .

이 방법은 단순히 모델 내부 분산만 사용하는 기존 접근과 달리 , 실제 예측 오차 분포를 반영하여 금융 리스크 관리에 더 적합한 예측구간을 제공한다는 장점이 있다 [19].

또한 구현 측면에서 분위수 (quantile) 구간으로도 예측구간을 제공한다 ( 예 : 10\% – 90\% 분위 ). forecast.py 기본 설정은 , 를 사용한다 .

\subsubsection{국가 간 환율 격차 정의}
국가 i 와 j 의 상대적 환율 갭을 예측 평균 기반으로 정의한다 .

이는 국가 간 상대적 강세 / 약세를 정량화하며 , 분위수 구간을 함께 제시하면 갭의 불확실성까지 해석할 수 있다 . 예측구간을 함께 고려하면 다음과 같이 표현된다 .

\section{에이전트 기반 모델 검증}
\subsection{에이전트 시스템 설계 (Framework Figure 생성 )}
외환 시장을 둘러싼 복잡하고 상호 연결된 거시경제적 (Macroeconomic) 요인들을 체계적으로 분석하기 위해 , 본 연구는 LangGraph 프레임워크를 기반으로 한 다중 에이전트 워크플로우 (Multi-Agent Workflow) 를 제안한다 . 기존의 단일 LLM(Large Language Model) 에 의존하는 프롬프팅 (Prompting) 방식과 달리 , 본 시스템은 그래프 기반의 상태 머신 (State Machine) 구조 (Derouiche et al., 2025)[] 를 채택하였다 . 이를 통해 각 분석 단계에서 도출된 에이전트의 중간 결과물 (State) 이 보존되며 , 이는 다음 단계 에이전트의 컨텍스트 (Context) 로 순차적으로 전달된다 . 이러한 설계 구조는 소형 언어 모델 (SLM, Liu et al., 2026)[] 이 지닌 연산 리소스의 한계를 극복하고 , 제한된 맥락 (Context Window) 내에서도 복잡한 경제적 추론의 일관성과 심도를 확보하는데 기여한다 .

\subsection{그래프 기반 RAG}
거시경제 데이터와 환율의 관계성 사이 선형적 인과관계를 효과적으로 검증하기 위해 , 본 시스템은 LightRAG(Guo et al., 2024)[] 를 핵심 지식 검색 증강 (eval Augmented Generation) 엔진으로 통합하였다 . 기존의 단순 벡터 속성 기반 검색과 달리 , 제안하는 방식은 IMF, OECD, BIS 등 기관의 실제 경제 보고서에서 추출한 내용 (Entity) 와 그들 사이의 인과성 (Relationship) 을 지식 그래프 (Knowledge Graph) 형태로 인덱싱 (Indexing) 한다 . 이는 단편적인 텍스트 매칭을 넘어 , 그래프 구조로 에이전트가 보다 고차원적이고 맥락 기반의 펀더멘털 (Fundamental) 분석을 수행할 수 있도록 지식의 기반을 제공한다 .

\subsection{다중 에이전트 역할 정의}
본 시스템의 파이프라인은 실제 경제 분석 프로세스를 모사하여 5 개의 전문화된 에이전트들로 구성되었다 .

\begingroup\small
\begin{table}[t]
\centering
\caption{Table n. 에이전트 역할 정의}
\begin{tabular}{|p{0.25\linewidth}|p{0.70\linewidth}|}
\hline
 역할 & 내용 \\
\hline
 계량경제 분석가 & 모델의 예측 결과의 통계적 유의성 (Statistical significance) 및 분산 검증 \\
\hline
 거시경제 분석가 & 경제 성장률 , 인플레이션 등 거시경제의 펀더멘털 (Fundamental) 요인 분석 \\
\hline
 통화정책 전략가 & 중앙은행의 금리 , 가계부채 등 통화기조가 환율 변화에 미치는 영향 모델링 \\
\hline
 테일 리스트 평가관 & 관세 무역 , 정책 변화 등 모델이 포착하지 못하는 비정형적 요소 분석 및 평가 \\
\hline
 책임 경제학자 & 전 단계의 결과를 정량적 / 정석적 분석을 통하여 보고서형태의 결과 도출 \\
\hline
\end{tabular}
\end{table}
\endgroup

\subsection{에이전트 결과 (Pilot 결과 참고하여 Appendix 추가 )}
제안하는 다중 에이전트 시스템의 결과를 검증한 결과 , 본 모델은 과거 시계열 데이터의 추세나 연장에만 의존하는 기존 정량모델 ( 비교 모델 언급 ) 의 근본적인 한계를 극복하였다 . 특정 경제 이벤트 발생 시 해당 시스템은 모델의 예측 값과 신뢰구간 탐색으로 불확실성을 식별하고 , 각 전문 에이전트의 교차 검증을 통해 구조적 변화와 거시경제적 요인의 영향을 성공적으로 검증하고 탐색하였다 . 이는 본 시스템이 단순한 텍스트 생성이 아닌 , 금융 시장의 복합적 요소를 심층적으로 분석하고 해석할 수 있는 타당성을 지니고 있음을 입증한다 .

\section{결과론}
본 절에서는 제안한 Bayesian Multivariate Time-series Graph Neural Network(B‑MTGNN) 을 글로벌 외환 (FX) 시장의 국가 단위 환율 예측에 적용한 실증 결과를 제시한다 . 분석 대상은 미국 (US Trade Weighted Dollar Index), 한국 (KRW/USD), 일본 (JPY/USD) 3 개 통화와 33 개 거시경제 변수로 구성된 월별 다변량 시계열이며 , 국가 간 상호의존성과 비정상성이 강한 환경에서 네트워크 기반 · 불확실성 인지형 예측의 성능을 평가하는 데 목적이 있다 .

\subsection{연구 설정 및 평가 설계}
실험에는 2011 년 1 월 \textasciitilde{}2025 년 12 월 (180 개월 ) 의 월별 데이터가 사용되었으며 , 학습 구간은 2011 년 1 월 \textasciitilde{}2023 년 12 월 (156 개월 ), 검증 구간은 2024 년 1 월 \textasciitilde{}12 월 (12 개월 ), 평가 구간은 2025 년 1 월 \textasciitilde{}12 월 (12 개월 ) 로 분할하였다 . 입력 시퀀스 길이는 , 출력 시퀀스 길이는 의 슬라이딩 윈도우 구조를 사용하여 단일 스텝 예측을 반복 연결함으로써 12 개월 (1 년 ) horizon 의 환율 경로를 생성하였다 .

주요 하이퍼파라미터는 학습률 0.00015, Dropout 0.02, 은닉 차원 256, 학습 에폭 180 이며 , Bayesian 불확실성 추정을 위해 MC dropout 전방추론 샘플 수를 학습 / 평가 시 , 장기 전망 (2026 년 예측 ) 시 으로 설정하였다 . 예측 성능 평가는 Root Relative Squared Error(RSE) 를 사용하였으며 , RSE < 0.5 를 목표 임계값으로 두어 성능을 판단하였다 .

모델 범용 안정성을 검증하기 위해 테스트 구간을 기준으로 과거 구간을 슬라이딩하는 롤링 백테스트를 수행하였다 . 180 개월에서 입력 길이 24 개월과 출력 길이 1 개월을 제외한 155 개 윈도우 중 145 개 윈도우에 대해 RSE 를 산출하고 평가하였다 . 이러한 설계는 비정상적 거시 / 금융 환경에서 예측 정확도와 불확실성 보정 수준을 동시에 검증하는 것을 목표로 한다 .

\subsection{불확실성 인식 예측 검증}
각 국가에 대해 MC dropout 을 활용하여 다수의 전방추론 샘플을 생성하고 , 이를 통해 예측 평균과 신뢰구간을 포함하는 다단계 환율 경로를 구성하였다 . 그림 n 은 관측 환율 ( 실선 ), 예측 평균 ( 점선 ), 그리고 예측 구간 ( 음영 ) 을 동시에 제시하며 , 예측 시계가 길어질수록 신뢰구간이 점진적으로 확장되어 장기 FX 예측에 내재된 인식론적 불확실성 증가를 반영함을 보여준다 .

테스트 구간의 점추정 성능은 RSE 기준으로 평가되었으며 , 세 타겟 모두에서 RSE < 0.5 로 나타났다 . 이는 예측 오차가 실제 변동성 대비 상대적으로 낮은 수준임을 보여준다 .

\begingroup\small
\begin{table}[t]
\centering
\caption{표 2. 대상 국가별 RSE 성능 평가 테스트}
\begin{tabular}{|p{0.25\linewidth}|p{0.70\linewidth}|}
\hline
 국가 & RSE \\
\hline
 US (Trade Weighted Dollar Index) & 0.3921 \\
\hline
 JP (jp\_fx) & 0.2671 \\
\hline
 KR (kr\_fx) & 0.3570 \\
\hline
\end{tabular}
\end{table}
\endgroup

\subsection{롤링 안정성 분석}
시계열 전반의 예측 성능을 평가하기 위해 , 우리 연구에서는 테스트 구간 전체에 대해 슬라이딩 윈도우 평가를 수행했다 . Rolling RSE 는 입력 시퀀스 길이 일 때 , 모델은 매 시점마다 입력 윈도우를 갱신하면서 1-step 예측을 순차적으로 생성한다 . 각 윈도우 위치마다 예측 오차를 독립적으로 계산하여 , 전체 구간에 걸친 성ㄴ으을 평가할 수 있도록 했다 .

Relative Squared Error(RSE) 는 다음과 같이 정의한다 .

\subsection{잔차 기반 예측 간격 및 불확실성 분해}
Dropout 비율 0.02 에서 MC dropout 만으로 구성된 예측구간은 폭이 지나치게 좁아 실제 오차를 충분히 포괄하지 못하는 한계가 있었다 . 이는 MC dropout 이 주로 모델 파라미터의 인식론적 불확실성만을 포착하고 , 관측 잡음 등 우연적 불확실성은 직접 반영하지 못하기 때문이다 .

이를 보완하기 위해 검증 구간 잔차 분산 을 MC dropout 분산 에 결합하고 , 조정 계수 를 적용한 잔차 기반 예측구간을 도입하였다 . 이 결합 분산 접근은 불확실성을 동시에 반영해 기존 단일 원천 기반 구간보다 보정과 실용성이 개선된 구간을 제공한다 . jp\_fx 2025 년 테스트 결과에서 , 결합 구간은 실제 경로를 대부분 포괄하면서도 과도하게 넓지 않아 리스크 평가 도구로 활용 가능한 균형을 보여주었다 .

\subsection{국가 간 환율 격차 역학}
국가 간 상대적 환율 격차는 예측 평균을 기반으로 로 정의하였다 . 해당 지표를 예측 horizon 에 따라 국가 , 상별로 상이한 동학적 패턴이 관측되었다 .

일부 쌍에서는 horizon 이 증가함에 다라 격차가 점진적으로 확대되는 단조 증가 패턴이 나타나는 반면 , 다른 쌍에서는 수렴 또는 진동적 패턴이 관측되어 정책 · 시장 조정 메커니즘이 작동하고 있음을 관측하였다 . 한국 원화 – 일본 엔화 조합은 예측 구간 전반에 걸쳐 상대적으로 동적인 격차 변동을 확인하였으며 , 격차 기울기 분석 결과 , 특정 horizon 구간에서 지속적인 양의 기울기가 관측되었다 .

또한 , 격차에 대한 예측 불확실성 구간을 함께 분석한 결과 , 일부 국가 쌍에서는 horizon 확장에 따라 신뢰구간 폭이 점진적으로 확대되는 경향이 나타났으며 , 이는 장기 예측 구간에서 상대적 격차의 분산이 증가함을 보여준다 .

전반적으로 국가 쌍별 경로는 동일한 방향으로 움직이지 않았으며 , 각 조합마다 상이한 변동성 수준과 조정 속도가 관측되었다 .

\subsection{2026 년 전망 및 연구 결과}
마지막으로 2011–2025 년 데이터를 모두 활용하여 2026 년 1 월 \textasciitilde{}12 월에 대한 환율 전망을 예측하였다 . 이는 2025 년 12 월까지의 관측값을 조건으로 MC dropout 을 사용해 생성한 확률적 예측을 요약한 결과이며 , 표 4 는 예측 평균과 95\% 신뢰구간 , CI 폭 / 평균 비율을 제시한다 .

\begingroup\small
\begin{table*}[t]
\centering
\caption{표 4. 2026 년 환율 예측 요약}
\begin{tabular}{|p{0.14\linewidth}|p{0.14\linewidth}|p{0.14\linewidth}|p{0.14\linewidth}|p{0.36\linewidth}|}
\hline
 국가 & 예측 평균 & 95\% CI 하한 & 95\% CI 상한 & CI 폭 / 평균 비율 \\
\hline
 US (TWDI) & 118.45 & 118.21 & 118.69 & ±0.20\% \\
\hline
 KR (KRW/USD) & 1,451.19 & 1,448.09 & 1,454.29 & ±0.21\% \\
\hline
 JP (JPY/USD) & 149.67 & 149.52 & 149.82 & ±0.10\% \\
\hline
\end{tabular}
\end{table*}
\endgroup

MC dropout 기반 신뢰구간은 평균 대비 ±0.1\textasciitilde{}0.2\% 로 매우 좁게 형성되어 , Dropout 0.02 설정 하에서는 인식론적 분산만 제한적으로 반영된 결과로 해석된다 .

국가 간 환율 격차 의 2026 년 전망에서는 한국 – 일본 간 격차가 가장 높은 동학성과 불확실성을 동반하는 것으로 나타났고 , 분위수 구간 제시를 통해 잠재적 발산 · 수렴 경로와 그 불확실성 범위를 동시에 평가할 수 있었다 . 다만 2026 년은 아직 관측되지 않은 미래 구간이므로 , 본 전망은 2025 년 테스트에서 확보된 RSE 0.27\textasciitilde{}0.49 수준의 성능에 기반한 방향 · 추세 수준의 참고 지표로 해석해야 하며 , 절대 수준에 대해서는 향후 실측치에 따른 지속적 모니터링과 재보정이 필요하다 .

이러한 결과는 B‑MTGNN 이 단순 점추정 정확도를 넘어 , 네트워크 기반 통화 동학과 불확실성을 통합적으로 반영하여 국제 포트폴리오 배분 · 환율 위험 관리 · 정책 분석에 활용 가능한 의사결정 지향적 분석 프레임워크를 제공함을 시사한다 .

\section{Discussion}
\subsection{비대칭적 통화 공진화의 구조적 성격}
본 연구는 국가 간 환율 동학이 일관된 비대칭성을 보인다는 점을 확인하였다 . 환율 격차 (exchange rate gap) 분석에서 국가 쌍별로 발산 (divergence) 및 수렴 (convergence) 패턴이 상이하게 나타났으며 , 이러한 차이는 단일 국가의 거시경제 여건이나 개별 양자 관계만으로는 충분히 설명되기 어렵다 . 이는 무역 연계 , 금융 통합 , 지역 파급효과로 구성된 상호작용 네트워크에서 각 국가가 점유하는 구조적 위치가 상대적 통화 진화 경로에 체계적으로 반영됨을 시사한다 .

나아가 비대칭적 공진화는 일시적 통계 현상이 아니라 글로벌 FX 시스템의 구조적 속성으로 해석될 수 있다 . 네트워크 중심성이 높은 국가는 인접 통화의 동학에 상대적으로 큰 영향을 미치는 경향을 보이며 , 이는 지속적인 상대적 절상 / 절하 압력으로 연결될 수 있다 . 이러한 결과는 FX 시장을 복잡적응시스템 (complex adaptive system) 으로 보는 관점과 부합하며 , 국지적 상호작용이 시스템 수준의 창발적 (emergent) 거동을 유도할 수 있음을 뒷받침한다 .

\subsection{네트워크 구조와 시스템적 결합의 역할}
네트워크 수준의 분석 결과 , 국가 간 환율 동학은 글로벌 FX 시장 전반에서 이질적으로 분포하였다 . 경제적 · 금융적으로 밀접한 국가 집단에서는 통화 움직임의 동조화가 강화되고 충격 전파 속도가 빨라지는 높은 시스템적 결합 (systemic coupling) 영역이 관찰되었다 . 반면 , 결합도가 낮은 영역에서는 환율 변동이 비교적 독립적이며 개별적 (idiosyncratic) 특성을 보였다 .

특히 시스템적 결합도가 상승하는 구간은 예측 불확실성의 증가와 동반되는 경향이 있었다 . 이는 글로벌 금융 스트레스 , 정책 불확실성 , 유동성 재조정 (liquidity re-pricing) 국면에서 국가 간 상호의존성이 증폭될 수 있음을 의미한다 . 따라서 고립적 또는 양자 중심 예측 모형은 이러한 체제에서 위험을 과소평가할 가능성이 있으며 , 네트워크 구조를 명시적으로 반영하는 접근이 필요하다 .

\subsection{환율 예측 및 위험 분석에 대한 시사점}
본 연구는 환율 예측 및 거시금융 위험 분석에 다음의 시사점을 제공한다 .

첫째 , 절대적인 환율 경로의 정확도뿐 아니라 상대적 통화 동학을 함께 모델링하는 것이 중요하다 . 환율 격차의 동학은 통화 간 상대적 위치와 조정 경로가 체계적으로 진화할 수 있음을 보여준다 .

둘째 , 불확실성 인지형 접근은 구조적으로 유도된 발산과 단순 잡음에 의한 변동을 구분하는 원칙적 (principled) 수단을 제공한다 . 이는 장기 예측에서 특히 중요한 인식적 불확실성을 보다 신중하게 해석하게 하며 , 시나리오 분석 및 스트레스 테스트와 같은 전향적 위험 평가에 유용하다 .

셋째 , 시스템적 결합도 (systemic coupling) 및 중심성 (centrality) 과 같은 네트워크 기반 지표는 변동성 또는 양자 상관계수 기반의 전통적 위험 지표를 보완할 수 있다 . 본 연구의 프레임워크는 직접적인 매매 · 정책 최적화를 목표로 하지는 않으나 , 고결합 체제에서 분산 효과가 약화되고 네트워크 효과를 통해 통화 위험이 증폭될 수 있는 메커니즘에 대한 분석적 근거를 제공한다 .

\subsection{방법론적 기여 및 일반화 가능성}
방법론적으로 본 연구는 거시금융 시스템에 대한 그래프 기반 확률적 모델링 문헌에 기여한다 . 국가 간 의존성과 시간 동학을 동시에 학습하면서 인식적 불확실성을 정량화함으로써 , 제안한 B-MTGNN 은 결정론적 또는 점추정 중심의 기존 그래프 시계열 예측 모형의 한계를 보완한다 .

실증 분석은 제한된 주요 통화 집합을 대상으로 수행되었으나 , 제안 프레임워크는 개념적으로 확장 가능하다 . 비대칭적 공진화와 네트워크 상호작용이 두드러지는 다른 FX 바스켓 , 지역 블록 , 신흥시장 통화 집단에도 적용될 수 있으며 , 더 나아가 네트워크 상호작용과 불확실성 하 의사결정이 중요한 다른 복잡 시스템에도 적용 가능하다 .

\section{Conclusion}
본 연구는 글로벌 외환 (FX) 시장에 내재된 국가 간 통화의 비대칭적 공진화 (asymmetric co-evloutuion) 현상을 시스템 차원에서 규명하기 위해 , 새로운 딥러닝 기반 예측 프레임워크인 B-MTGNN(Bayesian Multivariate Time-series Graph Neural Network) 을 제안하고 그 유효성을 검증하였다 . 기존의 자산 중심적 (asset-centric) 이거나 양자 간 (bilateral) 비교에 국한되었던 환율 모델의 한계를 극복하고자 , 본 논문은 다변량 시계열 처리 알고리즘과 그래프 신경망 (GNN) 을 결합하여 통화 간의 상호작용을 모델링하였다 .

실증 분석 결과 , 개별 국가의 환율 동학은 구조적인 비대칭성을 띠며 , 네트워크 상호작용에 강하게 의존하여 진화함이 확인되었다 . 또한 , 제안된 프레임워크는 다양한 예측 시계 (forecasting horizon) 에 걸쳐 안정적으로 보정된 확률적 예측을 산출함을 입증하였다 . 특히 , 국가 간 환율 격차 분석을 통해 기존의 단순 예측 오차 지표만으로는 포착하기 어려운 통화 간의 지속적 발산과 수렴 , 그리고 체제 의존적 (regime-dependent) 동학을 심층적으로 규명하였다 . 네트워크 수준 분석은 국가별 이질적 동학과 글로벌 시스템적 결합이 동시에 존재함을 강조하며 , 무역 연계 , 거시경제 상호의존성 , 지역적 파급효과에 기인한 글로벌 환율 시스템 내의 파급효과를 포착하여 상대적 통화 움직임에 대한 통찰을 제공한다 .

본 연구는 단순한 매매 알고리즘이나 단기적 수익률 최적화 모델을 구축하는 것을 목적으로 하지 않는다 . 대신 , 고도의 불확실성이 지배하는 금융 복잡계 (complex systems) 를 해석하기 위한 분석 프레임워크를 제안한다는 점에서 학술적인 의의를 가진다 . 연구 결과는 국제 금융시장에서 상대적 환율의 진화 , 전염 효과 (contagion), 시스템적 위험을 분석할 때 네트워크 구조와 명시적 불확실성 정량화를 함께 고려하는 것이 중요함을 강조한다 . 더 나아가 , 제안된 접근법은 향후 이질적인 통화 집단 분석을 넘어 , 비대칭적 상호작용을 특징으로 하는 다양한 금융 도메인의 시계열 예측 문제로 확장 가능한 범용적 모델링 패러다임을 제공할 것으로 기대된다 .

참 고 문 헌 ( 중고딕 11pt, 가운데 맞춤 )

[1] S. Miranda-Agrippino and H. Rey, “US Monetary Policy and the Global Financial Cycle,” The Review of Economic Studies , vol. 87, no. 6, pp. 2754-2776, 2020.

[2] J. Baruník and T. Křehlík , “Measuring the Frequency Dynamics of Financial Connectedness and Systemic Risk,” Journal of Financial Econometrics , vol. 16, no. 2, pp. 271-296, 2018.

[3] Z. Pan, Y. Liang, W. Wang, Y. Yu, M. Li, and J. Zhang, “Urban Traffic Prediction from Spatio -Temporal Data: A Deep Learning Approach,” IEEE Transactions on Knowledge and Data Engineering , vol. 32, no. 4, pp. 720-733, 2019.

[4] D. Diebold and K. Yilmaz, “Better to Give than to Receive: Predictive Directional Measurement of Volatility Spillovers,” International Journal of Forecasting , vol. 28, no. 1, pp. 57-66, 2012.

[5 ] Wu , Z., Pan, S., Chen, G., Long, G., Zhang, C., \& Philip, S. Y. (2020). "Connecting the dots: Multivariate time series forecasting with graph neural networks." Proceedings of the 26th ACM SIGKDD International Conference on Knowledge Discovery \& Data Mining .

[6 ] Li , Y., Yu , R., Shahabi , C., \& Liu , Y. (2018). "Diffusion convolutional recurrent neural network: Data-driven traffic forecasting." International Conference on Learning Representations (ICLR) .

[7] Yu, S., et al. (2024). " FinAgent : A Multi-modal Generalist Agent for Financial Analysis and Reasoning." arXiv preprint arXiv:2402.14411 .

[8 ] Wu , Q., et al. (2023). " AutoGen : Enabling Next-Gen LLM Applications via Multi-Agent Conversation." Proceedings of the 40th International Conference on Machine Learning (ICML) .

[9] Pradeepkumar, D., \& Ravi, V. (2018). Forecasting financial time series volatility using particle swarm optimization trained quantile regression neural network. Applied Soft Computing , 58, 35-52.

[10] Fischer, T., \& Krauss, C. (2018). Deep learning with long short-term memory networks for financial market predictions.

[11] Zhou, H., Zhang, S., Peng, J., Zhang, S., Li, J., Xiong, H., \& Zhang, W. (2021). Informer: Beyond efficient transformer for long sequence time-series forecasting.

[12] Amat, C., Michalski, T., \& Stoltz, G. (2018). Fundamentals and exchange rate forecastability with machine learning methods. Journal of International Money and Finance , 88, 58-74.

[13] Yao, Q., Song, D., Chen, H., Wei, A., Cottrell, G. W., \& Bian, G. W. (2018). A Dual-Stage Attention-Based Recurrent Neural Network for Time Series Prediction.

[14] Yu, B., et al. (2018). Spatio -temporal graph convolutional networks: A deep learning framework for traffic forecasting. IJCAI .

[15] Kipf, T. N., \& Welling, M. (2017). Semi-supervised classification with graph convolutional networks. ICLR .

[16] Wu, Z., Pan, S., Long, G., Jiang, J., \& Zhang, C. (2019). Graph WaveNet for deep spatial-temporal graph modeling. Proceedings of the 28th International Joint Conference on Artificial Intelligence (IJCAI) , 1907-1913.

[17] Veli č kovi ć , P., Cucurull , G., Casanova, A., Romero, A., Lio, P., \& Bengio, Y. (2018). Graph Attention Networks. International Conference on Learning Representations (ICLR) .

18] Kendall, A., \& Gal, Y. (2017). “ What Uncertainties Do We Need in Bayesian Deep Learning for Computer Vision? ” Advances in Neural Information Processing Systems ( NeurIPS ), vol. 30, pp. 5574-5584.

[19] Gneiting, T., \& Raftery, A. E. (2007). “ Strictly Proper Scoring Rules, Prediction, and Estimation, ” Journal of the American Statistical Association, vol. 102, no. 477, pp. 359-378.

[3] C. W. Xu, “ Fuzzy Model Identification and Self-learning for Dynamic Systems, ” IEEE Trans. Syst. Man. Cybern . , vol. 17, no. 4, pp. 683-689, 1987.

[4] H. Takagi and M. Sugeno , “ Fuzzy Identification of Systems and Its Application to Modeling and Control, ” IEEE Trans. on Sys. Man and Cybern . , vol. 15, pp. 116-132, 1985.

[5] J. H. Holland, Genetic Algorithms in Search, Optimization and Machine Learning , Addison- Wesley, 1989. ( 신명조 9pt)

\end{document}